% Part: first-order-logic
% Chapter: proof-systems
% Section: natural-deduction

\documentclass[../../../include/open-logic-section]{subfiles}

\begin{document}

\iftag{FOL}
      {\olfileid{fol}{prf}{ntd}}
      {\olfileid{pl}{prf}{ntd}}

\olsection{Natuurlijke deductie}

Natuurlijke deductie is een stelsel van afleidingsregels dat echte
redeneringen probeert weer te geven. Het gaat vooral om redeneringen
van wiskundigen die precies volgens vaste regels verlopen. Zulke
redeneringen volgen een aantal vertrouwde patronen. Bij een bewijs door
gevalsonderscheid beginnen we met een premisse die uit twee delen met
``of'' bestaat. We laten zien dat de conclusie uit elk van beide delen,
de disjuncten, volgt. Bij een bewijs uit het ongerijmde laten we zien
dat de negatie van de conclusie tot een tegenspraak leidt. Om een
uitspraak met de vorm ``als \dots, dan \dots'' te bewijzen, nemen we
het deel na ``als'' aan en leiden we daaruit het deel na ``dan'' af.
Natuurlijke deductie legt enkele van deze vertrouwde manieren van
redeneren in formele regels vast. Bij elk logisch verbindingsteken en
elke kwantor horen twee regels: een introductieregel en een
eliminatieregel. Elke regel geeft een van deze manieren van
redeneren weer. Zo hoort $\Intro{\lif}$ bij het bewijzen van een
implicatie en $\Elim{\lor}$ bij een bewijs door gevalsonderscheid. Een
bijzonder eenvoudige regel is $\Elim{\land}$. Met die regel mogen we
uit $!A \land !B$ de formule~$!A$ afleiden, of de formule~$!B$.

Natuurlijke deductie gebruikt aannames en verschilt daarin van andere
stelsels van afleidingsregels. !!^a{derivation} in natuurlijke deductie
is een boom van !!{formula}s. Onderaan, aan de wortel van de boom,
staat één !!{formula}. De ``bladeren'' bovenaan zijn de !!{formula}s
waaruit de conclusie wordt afgeleid. Sommige formules aan de bladeren
zijn alleen tijdelijk beschikbaar. Ze spelen binnen de !!{derivation}
een rol, maar zijn niet meer beschikbaar wanneer we de conclusie
bereiken. Dat sluit aan bij gewone redeneringen waarin we een hypothese
voor korte tijd aannemen. Bij een bewijs door gevalsonderscheid nemen
we elk van de disjuncten apart aan. Bij het bewijzen van een
implicatie nemen we het deel na ``als'' aan. Bij een bewijs uit het
ongerijmde nemen we de negatie van de conclusie aan. In natuurlijke
deductie mogen we zo'n tijdelijke aanname invoeren, er een tussenstap
mee afleiden en de aanname daarna weer opheffen. Juist dat is een
kenmerk van dit systeem. De !!{formula}s aan de bladeren van
!!a{derivation} heten aannames. Sommige afleidingsregels kunnen een
aanname ``!!{discharge}.'' Stel bijvoorbeeld dat we !!a{derivation}
van~$!B$ hebben uit een aantal aannames, waaronder~$!A$. Met
$\Intro{\lif}$ mogen we dan~$!A \lif !B$ afleiden en elke aanname van
de vorm~$!A$ !!{discharge}. We geven de afleidingsstap en de aannames
die daar worden opgeheven hetzelfde nummer. Zo houden we bij welke
aanname bij welke stap wordt opgeheven. De aannames die aan het eind
van de !!{derivation} !!{undischarged} zijn, zijn samen voldoende om
de conclusie waar te maken. De !!{derivation} laat dus zien dat de
conclusie uit die overgebleven aannames volgt.

Voor natuurlijke deductie definiëren we de relatie $\Gamma \Proves !A$
als volgt. Zij geldt precies als er !!a{derivation} bestaat waarin
$!A$~de laatste !!{sentence} van de boom is en elk blad dat
!!{undischarged} is tot~$\Gamma$ behoort. $!A$~is een stelling in
natuurlijke deductie precies als er !!a{derivation} bestaat waarin
$!A$~de laatste !!{sentence} is en alle aannames zijn !!{discharged}.
De volgende !!a{derivation} laat bijvoorbeeld zien dat $\Proves (!A
\land !B) \lif !A$:
\begin{prooftree}
  \AxiomC{$\Discharge{!A \land !B}{1}$}
  \RightLabel{\Elim{\land}}
  \UnaryInfC{$!A$}
  \DischargeRule{\Intro{\lif}}{1}
  \UnaryInfC{$(!A \land !B) \lif !A$}
\end{prooftree}
Het nummer~$1$ laat zien dat de aanname $!A \land !B$ bij de
\Intro{\lif}-stap wordt !!{discharged}.

Een verzameling~$\Gamma$ is inconsistent in natuurlijke deductie
precies als $\Gamma \Proves \lfalse$. Door de regel \FalseInt{} kan uit
een inconsistente verzameling elke !!{sentence} worden !!{derive}d.

Gerhard Gentzen en Stanis\l{}aw Ja\'skowski ontwikkelden in de jaren
dertig van de twintigste eeuw systemen voor natuurlijke deductie. Dag
Prawitz en Frederic Fitch werkten die systemen later verder uit. De
afleidingsstappen lijken op manieren waarop mensen van nature bewijzen.
Daarom gebruiken filosofen natuurlijke deductie graag. De versies van
Fitch komen vaak voor in inleidende leerboeken over logica. Binnen de
filosofie van de logica worden de regels van de natuurlijke deductie
soms opgevat als regels die de betekenis van de logische operatoren
bepalen. Deze opvatting heet ``bewijstheoretische semantiek''.
\end{document}
